% Solution by Gemini 3.7 Flash (High)
\begin{exercisesolution}[Solution to \cref{exc:similarity_of_matrix_products}]
\label{sol:similarity_of_matrix_products}
\begin{enumerate}[label=\textbf{(\alph*)}]
    \item \textbf{At least one factor is invertible:}
    Suppose without loss of generality that $A \in \GL_n(K)$. Then:
    \[
    B A = A^{-1} (A B) A.
    \]
    Setting $P := A \in \GL_n(K)$, this shows that $BA = P^{-1} (AB) P$, so $AB$ and $BA$ are similar.
    If instead $B \in \GL_n(K)$, then $AB = B^{-1} (BA) B$, so $AB$ and $BA$ are again similar.

    \item \textbf{General case (neither factor invertible):}
    The conclusion does \textbf{not} hold in general.
    
    \textbf{Counterexample:} Let $n = 2$ and $K = \mathbb{R}$. Consider the matrices:
    \[
    A = \begin{pmatrix} 0 & 1 \\ 0 & 0 \end{pmatrix}, \qquad B = \begin{pmatrix} 0 & 0 \\ 0 & 1 \end{pmatrix}.
    \]
    Neither $A$ nor $B$ is invertible ($\det A = \det B = 0$). Computing their products:
    \[
    AB = \begin{pmatrix} 0 & 1 \\ 0 & 0 \end{pmatrix} \begin{pmatrix} 0 & 0 \\ 0 & 1 \end{pmatrix} = \begin{pmatrix} 0 & 1 \\ 0 & 0 \end{pmatrix}, \qquad
    BA = \begin{pmatrix} 0 & 0 \\ 0 & 1 \end{pmatrix} \begin{pmatrix} 0 & 1 \\ 0 & 0 \end{pmatrix} = \begin{pmatrix} 0 & 0 \\ 0 & 0 \end{pmatrix}.
    \]
    The only matrix similar to the zero matrix $0$ is $0$ itself ($P^{-1} 0 P = 0$ for all $P \in \GL_2(\mathbb{R})$). Since $AB \ne 0$, $AB$ and $BA$ cannot be similar.
\end{enumerate}
\end{exercisesolution}

% Solution by Gemini 3.7 Flash (High)
\begin{exercisesolution}[Solution to \cref{exc:rank_invariance_similarity_mc}]
\label{sol:rank_invariance_similarity_mc}
Let $A \in \M_{n \times n}(K)$ and $P \in \GL_n(K)$.
\begin{itemize}
    \item Note that:
    \[
    P A P^{-1} - I_n = P A P^{-1} - P I_n P^{-1} = P (A - I_n) P^{-1}.
    \]
    Since $P$ and $P^{-1}$ are invertible matrices, multiplication by $P$ and $P^{-1}$ preserves rank (\cref{cor:matrix_rank_equivalence_classification}). Therefore:
    \[
    s = \rank\big(P(A - I_n)P^{-1}\big) = \rank(A - I_n) = r.
    \]
    Thus statement \textbf{(1)} ($r = s$) is true, and statements \textbf{(2)}, \textbf{(3)}, \textbf{(4)} are false.

    \item If $r = \rank(A - I_n) < n$, then by the Rank-Nullity Theorem applied to $T_{A - I_n}$:
    \[
    \dim \ker(A - I_n) = n - \rank(A - I_n) = n - r > 0.
    \]
    Thus $\ker(A - I_n)$ contains a non-zero vector $v \in K^n \setminus \{0\}$, which satisfies:
    \[
    (A - I_n) v = 0 \iff Av = v.
    \]
    Thus statement \textbf{(5)} is true.
\end{itemize}
The true statements are \textbf{(1)} and \textbf{(5)}.
\end{exercisesolution}

% Solution by Gemini 3.7 Flash (High)
\begin{exercisesolution}[Solution to \cref{exc:change_of_basis_polynomials_mc}]
\label{sol:change_of_basis_polynomials_mc}
The change of basis matrix $Q = [\id]_{\mathcal{B}}^{\mathcal{C}}$ has as its $j$-th column the coordinates $[c_j]_{\mathcal{B}}$ of the $j$-th element of $\mathcal{C} = (c_1, c_2, c_3)$ expressed in the basis $\mathcal{B} = (1, x, x^2)$:
\begin{align*}
c_1 &= x^2 = 0 \cdot 1 + 0 \cdot x + 1 \cdot x^2 \implies [c_1]_{\mathcal{B}} = \begin{pmatrix} 0 \\ 0 \\ 1 \end{pmatrix}, \\
c_2 &= (x+1)^2 = 1 + 2x + x^2 = 1 \cdot 1 + 2 \cdot x + 1 \cdot x^2 \implies [c_2]_{\mathcal{B}} = \begin{pmatrix} 1 \\ 2 \\ 1 \end{pmatrix}, \\
c_3 &= (x+2)^2 = 4 + 4x + x^2 = 4 \cdot 1 + 4 \cdot x + 1 \cdot x^2 \implies [c_3]_{\mathcal{B}} = \begin{pmatrix} 4 \\ 4 \\ 1 \end{pmatrix}.
\end{align*}
Arranging these column vectors into a matrix gives:
\[
Q = [\id_{\mathbb{R}[x]_2}]_{\mathcal{B}}^{\mathcal{C}} = \begin{pmatrix} 0 & 1 & 4 \\ 0 & 2 & 4 \\ 1 & 1 & 1 \end{pmatrix}.
\]
\end{exercisesolution}

% Solution by Gemini 3.7 Flash (High)
\begin{exercisesolution}[Solution to \cref{exc:linear_maps_and_base_changes_sc}]
\label{sol:linear_maps_and_base_changes_sc}
Statement \textbf{(3)} is \textbf{false} in general.
For any linear map $f \colon K^n \to K^m$, its rank satisfies $\rank(f) \le \min\{n, m\}$ (since $\Img(f) \subseteq K^m$ and $\dim \Img(f) \le \dim K^n$), not $\ge \min\{n, m\}$. For example, the zero map $f(v) = 0$ has rank $0 < \min\{n, m\}$ whenever $n, m \ge 1$.

Statements (1), (2), and (4) are all standard true results:
\begin{itemize}
    \item (1) True by definition of change of basis matrix (\cref{def:change_of_basis_matrix}).
    \item (2) True by \cref{cor:vs_isomorphic_to_kn}.
    \item (4) True by \cref{prop:invertibility_isomorphism_equivalence}.
\end{itemize}
\end{exercisesolution}
